\FifthArticle{11}
Zur numerischen Bestimmung der Wurzeln algebraischer Gleichungen, die nur aus drei Gliedern bestehen, lassen sich verschiedene Methoden anwenden. Sie können hier eine Eleganz und Bequemlichkeit erreichen, hinter denen die mühsamen Operationen weit zurückbleiben, die bei Gleichungen weniger einfacher Gestalt unvermeidlich sind. Solche Methoden verdienen daher wohl eine eigene Darstellung, zumal Gleichungen dieser Form häufig genug vorkommen.

Dies gilt zunächst für die Entwicklung der Wurzeln in unendliche Reihen. Tatsächlich lässt sich \textit{jede} Wurzel einer Gleichung mit drei Gliedern, gleichgültig ob reell oder imaginär, durch eine konvergente Reihe mit einfachem Bildungsgesetz ausdrücken. Ich werde jedoch \textit{diese} Lösungsart aus mehreren Gründen von meiner gegenwärtigen Betrachtung ganz ausschließen. Hier bemerke ich nur, dass die Konvergenzgeschwindigkeit vom Verhältnis der Koeffizienten zueinander abhängt: Sie ist umso langsamer, je näher dieses Verhältnis demjenigen kommt, bei dem die Gleichung zwei gleiche Wurzeln hat, und in diesem Grenzfall selbst ist sie langsamer als bei jeder fallenden geometrischen Folge. So bemerkenswert diese Reihen in allgemeiner theoretischer Hinsicht auch sind, wird man in der Praxis doch immer den \textit{indirekten} Methoden den Vorzug geben, die in den folgenden Artikeln entwickelt werden sollen, abgesehen von dem Fall, dass die Reihen sehr schnell konvergieren.

\FifthArticle{12}
Um die \textit{reellen} Wurzeln zu finden, benutze ich meine erstmals 1810 gedruckte Hilfstafel für Logarithmen von Summen und Differenzen oder, wenn eine größere Genauigkeit verlangt wird, als fünfstellige Logarithmen bieten\SourceLeaf{093} können, die ähnliche, aber erweiterte Tafel von \textsc{Matthiessen}. Einige besondere Anwendungen dieses Verfahrens habe ich schon früher veröffentlicht: zur Lösung quadratischer Gleichungen bei der 1840 erschienenen zwanzigsten Ausgabe von \textsc{Vega}s logarithmischem Handbuch und zur Lösung der kubischen Gleichung, die bei parabolischer Bewegung zur Bestimmung der wahren Anomalie dient, in Nr. 474 der Astronomischen Nachrichten. Dort wurde auch schon auf die allgemeine Anwendbarkeit des Verfahrens auf alle algebraischen Gleichungen mit drei Gliedern hingewiesen. Obwohl die Ausführung dieses ganz elementaren Gegenstands keine Schwierigkeiten bereitet, wird man bei der recht großen Vielfalt der Fälle doch gern einige Seiten für ihre übersichtliche Einteilung und die Zusammenstellung gebrauchsfertiger Vorschriften eingeräumt sehen.

Anstelle jener logarithmischen Hilfstafeln kann man auch die gewöhnlichen logarithmisch-trigonometrischen Tafeln verwenden. Die ersteren sind jedoch im Allgemeinen für den vorliegenden Zweck bequemer und bieten eine doppelt so große Genauigkeit wie die letzteren. Ich würde daher den Gebrauch trigonometrischer Tafeln für diese Aufgabe auf den seltenen Fall beschränken, dass man auch die mit siebenstelligen Logarithmen erreichbare Genauigkeit noch übertreffen möchte und dafür die bekannten zehnstelligen Logarithmen in \textsc{Vlacq}s oder \textsc{Vega}s Thesaurus verwenden kann. Übrigens sind beim Gebrauch der Hilfslogarithmen doppelt so viele Fälle zu unterscheiden wie beim Gebrauch der trigonometrischen Logarithmen. Als Nachteil darf dies jedoch nicht gelten: Liegt einmal die vollständige allgemeine Einteilung vor, lässt sich jeder einzelne Fall leicht seiner Klasse zuordnen, und die eigentliche indirekte Aufgabe lässt sich umso leichter ausführen, wenn die ganze Klasse nur den halben Umfang hat. Gerade deshalb lässt sich aber die allgemeine Einteilung für die Lösung mit trigonometrischen Logarithmen kürzer und bequemer darstellen. Ich werde sie daher vorausschicken, da sich die Einteilung für die andere Lösungsform dann von selbst daraus ergibt.

\FifthArticle{13}
Die Methode wird unmittelbar nur zur Bestimmung der \textit{positiven} Wurzeln einer gegebenen Gleichung angewandt. Die negativen Wurzeln\SourceLeaf{094} erhält man, indem man dasselbe Verfahren auf die Gleichung anwendet, die aus jener durch Einführung der zur ursprünglichen Unbekannten entgegengesetzten Größe entsteht.

Ich setze die Gleichung in die Form
\[
x^{m+n}\pm e x^m\pm f=0
\]
wobei $m$, $n$, $e$, $f$ gegebene \textit{positive} Größen bedeuten. Diese Form umfasst je nach der Kombination der Vorzeichen eigentlich vier verschiedene Fälle. Der erste, bei dem beide Male die oberen Vorzeichen gewählt werden, fällt jedoch weg, da die Gleichung
\[
x^{m+n}+e x^m+f=0
\]
offenbar keine positive Wurzel haben kann. Übrigens darf man voraussetzen, dass $m$ und $n$ (worunter ganze Zahlen verstanden werden, obwohl die Anwendbarkeit der Methode an sich davon unabhängig ist) keinen gemeinsamen Teiler haben, da sich jeder andere Fall leicht auf diesen zurückführen lässt. Schließlich schreibe ich zur Abkürzung
\[
\frac{f^n}{e^{m+n}}=\lambda
\]

\Needspace{6\baselineskip}\begin{center}\textit{Erste Form.}\end{center}
\[
x^{m+n}+e x^m-f=0
\]
Man führe einen Winkel $\theta$ ein, der immer im ersten Quadranten zu nehmen ist, sodass
\[
\frac{x^{m+n}}{f}=\sin\theta^2,\quad
\frac{e x^m}{f}=\cos\theta^2
\]
und somit (I)
\[
x^{m+n}=f\sin\theta^2,\quad
x^m=\frac{f\cos\theta^2}{e},\quad
x^n=e\mathrm{tang}\theta^2
\]
gilt. Durch Elimination von $x$ ergibt sich die Gleichung
\[
\lambda=\frac{\sin\theta^{2m}}{\cos\theta^{2m+2n}}
\]
aus der $\theta$ bestimmt werden muss. Man erkennt leicht, dass die rechte Seite dieser Gleichung, als Funktion einer unbestimmten Größe $\theta$ betrachtet, von $0$\SourceLeaf{095} bis $\infty$ wächst, während $\theta$ alle Werte von $0$ bis $90^0$ durchläuft. Daher gibt es einen und nur einen Wert von $\theta$, der jener Gleichung genügt. Ist er gefunden, erhält man $x$ aus einer der Formeln I. Man beachte, dass $\theta=45^0$ für $\lambda=2^n$ wird und daher $\theta$ im ersten Oktanten zu suchen ist, wenn $\lambda$ kleiner, und im zweiten, wenn $\lambda$ größer als $2^n$ ist.

\Needspace{6\baselineskip}\begin{center}\textit{Zweite Form.}\end{center}
\[
x^{m+n}-ex^m-f=0
\]
Hier setzt man
\[
fx^{-m-n}=\sin\theta^2,\quad ex^{-n}=\cos\theta^2
\]
oder (I)
\[
x^{m+n}=\frac{f}{\sin\theta^2},\quad
x^n=\frac{e}{\cos\theta^2},\quad
x^m=\frac{f\mathrm{cotang}\theta^2}{e}
\]
Demnach ist $\theta$ aus der Gleichung
\[
\lambda=\frac{\sin\theta^{2n}}{\cos\theta^{2m+2n}}
\]
zu bestimmen, was auf eine und nur eine Weise geschehen kann. Der Wert von $x$ ergibt sich dann aus einer der Gleichungen I. $\theta$ liegt im ersten oder zweiten Oktanten, je nachdem, ob $\lambda$ kleiner oder größer als $2^m$ ist.

\Needspace{6\baselineskip}\begin{center}\textit{Dritte Form.}\end{center}
\[
x^{m+n}-ex^m+f=0
\]
Hier setzt man
\[
\frac{x^n}{e}=\sin\theta^2,\quad
\frac{fx^{-m}}{e}=\cos\theta^2
\]
oder (I)
\[
x^{m+n}=f\mathrm{tang}\theta^2,\quad
x^m=\frac{f}{e\cos\theta^2},\quad
x^n=e\sin\theta^2
\]
Eine dieser Formeln dient zur Bestimmung von $x$, sobald der Wert von $\theta$ gefunden ist. Dieser ergibt sich durch Lösung der Gleichung
\[
\lambda=\cos\theta^{2n}\sin\theta^{2m}
\]
Da die rechte Seite dieser Gleichung, als Funktion einer unbestimmten Größe $\theta$ betrachtet, sowohl für $\theta=0$ als auch für $\theta=90^0$ verschwindet, muss dazwischen ein größter Wert liegen. Da das Differential des\SourceLeaf{096} Logarithmus dieser Funktion $=(2m\mathrm{cotg}\theta-2n\mathrm{tang}\theta)\,\mathrm{d}\theta$ ist, tritt der größte Wert für $\theta=\theta^*$ ein, wenn man $\surd\frac{m}{n}=\mathrm{tang}\theta^*$ setzt. Die Funktion nimmt demnach von $0$ bis zu ihrem größten Wert, der offenbar
\[
=\frac{m^m n^n}{(m+n)^{m+n}}
\]
ist, zu und von dort bis $0$ ab, während $\theta$ von $0$ bis $\theta^*$ und von dort bis $90^0$ zunimmt. Der Maximalwert ist daher stets größer als der Wert für $\theta=45^0$, also größer als $\frac{1}{2^{m+n}}$, außer wenn $m=n$ und somit $\frac{1}{2^{m+n}}$ selbst der Maximalwert ist.

Daraus folgt, dass, je nachdem, ob $\lambda$ größer als
\[
\frac{m^m n^n}{(m+n)^{m+n}}
\]
oder kleiner ist, die Gleichung $\lambda=\cos\theta^{2n}\sin\theta^{2m}$ gar nicht oder durch zwei verschiedene Werte von $\theta$ erfüllt werden kann. Im ersten Fall hat die Gleichung $x^{m+n}-ex^m+f=0$ keine (positive) Wurzel, im anderen zwei. Im besonderen Fall, dass
\[
\lambda=\frac{m^m n^n}{(m+n)^{m+n}}
\]
gilt, fallen beide Lösungen zusammen, und die Gleichung hat zwei gleiche Wurzeln. Dafür kann man nach Belieben eine der drei Formeln verwenden:
\[
x^{m+n}=\frac{fm}{n},\quad
x^m=\frac{f(m+n)}{en},\quad
x^n=\frac{em}{m+n}
\]

Was die Oktanten betrifft, in die die Werte von $\theta$ fallen, wenn tatsächlich zwei Lösungen vorhanden sind, so erkennt man leicht: Ist $\lambda$ größer als $\frac{1}{2^{m+n}}$, so liegen beide Werte von $\theta$ mit $\theta^*$ im selben Oktanten, und zwar im ersten oder zweiten, je nachdem, ob $m$ kleiner oder größer als $n$ ist. Ist dagegen $\lambda$ kleiner als $\frac{1}{2^{m+n}}$, so ist ein Wert von $\theta$ im ersten, der andere im zweiten Oktanten zu suchen. Im besonderen Fall, dass $\lambda=\frac{1}{2^{m+n}}$, ist $45^0$ selbst ein Wert von $\theta$, und der andere liegt im selben Oktanten wie $\theta^*$.

Noch sei die aus dieser Untersuchung aller drei Formen leicht folgende Konsequenz bemerkt, dass unsere Gleichung (sofern wir annehmen, dass $m$ und $n$ keinen gemeinsamen Teiler haben) nicht mehr als drei reelle Wurzeln haben kann, was auch aus anderen Gründen bekannt ist.
\SourceLeaf{097}

